    \documentclass[12pt,a4paper,oneside]{report}

    % ==============================================================================
    % Packages & Formatting (DTU Academic Guidelines)
    % ==============================================================================
    \usepackage[utf8]{inputenc}
    \usepackage[margin=1in, top=1.2in, bottom=1.2in]{geometry}
    \usepackage{setspace}
    \onehalfspacing
    \usepackage{amsmath, amssymb, amsfonts}
    \usepackage{booktabs}
    \usepackage{multirow}
    \usepackage{graphicx}
    \usepackage{float}
    \usepackage{xcolor}
    \usepackage{listings}
    \usepackage{fancyhdr}
    \usepackage{titlesec}
    \usepackage[hidelinks]{hyperref}
    \usepackage{tabularx}
    \usepackage{cite}
    \usepackage{enumitem}
    \usepackage{array}

    % Header and Footer setup
    \pagestyle{fancy}
    \fancyhf{}
    \fancyhead[L]{\fontsize{9}{11}\selectfont DTU/CSE/BTP/GP159: EdgeAgent -- Autonomous On-Device Mobile Agent}
    \fancyhead[R]{\fontsize{9}{11}\selectfont Minor Project Detailed Report}
    \fancyfoot[C]{\thepage}
    \renewcommand{\headrulewidth}{0.4pt}
    \renewcommand{\footrulewidth}{0.4pt}

    % Section title formatting
    \titleformat{\chapter}[display]
    {\normalfont\huge\bfseries}{\chaptertitlename\ \thechapter}{16pt}{\Huge}
    \titlespacing*{\chapter}{0pt}{0pt}{20pt}

    % Code Listing Formatting
    \lstset{
        basicstyle=\ttfamily\footnotesize,
        breaklines=true,
        frame=single,
        backgroundcolor=\color{gray!8},
        keywordstyle=\color{blue}\bfseries,
        commentstyle=\color{green!50!black}\itshape,
        stringstyle=\color{red!70!black},
        numbers=left,
        numberstyle=\tiny\color{gray},
        stepnumber=1,
        numbersep=8pt,
        tabsize=2,
        showstringspaces=false
    }

    \begin{document}

    % ==============================================================================
    % TITLE PAGE
    % ==============================================================================
    \begin{titlepage}
    \begin{center}
        {\large \textbf{DELHI TECHNOLOGICAL UNIVERSITY}}\\[0.2cm]
        {\normalsize (Formerly Delhi College of Engineering)}\\[0.2cm]
        {\normalsize Shahbad Daulatpur, Bawana Road, Delhi-110042}\\[0.6cm]

        {\large \textbf{DEPARTMENT OF COMPUTER SCIENCE \& ENGINEERING}}\\[1.2cm]

        {\Large \textbf{B. TECH Minor PROJECT REPORT}}\\[0.3cm]
        {\normalsize SESSION 2026--2027}\\[0.6cm]

    \rule{\linewidth}{1.2pt}\\[0.4cm]
    {\LARGE \textbf{EdgeAgent: An Autonomous On-Device Agent Architecture}}\\[0.25cm]
    {\LARGE \textbf{for Privacy-First ARM-Based Task Execution and Mobile Communication}}\\[0.3cm]
    \rule{\linewidth}{1.2pt}\\[1.0cm]

        \textit{Submitted in partial fulfillment of the requirements for the award of the degree of}\\[0.2cm]
        {\large \textbf{BACHELOR OF TECHNOLOGY}}\\[0.2cm]
        {\normalsize in}\\[0.2cm]
        {\large \textbf{COMPUTER SCIENCE \& ENGINEERING}}\\[1.0cm]

        \begin{minipage}{0.48\textwidth}
            \begin{flushleft} \large
                \textbf{Submitted by:}\\[0.2cm]
                \textbf{Ekansh Agarwal} (23/CS/149)\\
                \textbf{Faqre Alam} (23/CS/150)\\
                \textbf{Guneet Toppo} (23/CS/160)\\[0.2cm]
                \textbf{Junior Team Member:}\\
                \textbf{Bhawishya} (24/CS/137, 3rd Year)\\[0.2cm]
                \textbf{Group ID:} DTU/CSE/BTP/GP159
            \end{flushleft}
        \end{minipage}
        \hfill
        \begin{minipage}{0.48\textwidth}
            \begin{flushright} \large
                \textbf{Under the Supervision of:}\\[1.2cm]
                \rule{5.5cm}{0.4pt}\\[0.15cm]
                \textbf{(Signature of Supervisor)}\\[0.2cm]
                \textbf{Ms. Gull Kaur}\\
                Assistant Professor\\
                Department of CSE\\
                Delhi Technological University
            \end{flushright}
        \end{minipage}

        \vfill
        {\large Delhi Technological University • November 2026}
    \end{center}
    \end{titlepage}

    % ==============================================================================
    % CANDIDATE DECLARATION & CERTIFICATE
    % ==============================================================================
    \chapter*{Candidate's Declaration}
    \addcontentsline{toc}{chapter}{Candidate's Declaration}

In other words, we hereby certify that the work introduced in this Minor Project Report entitled \textbf{``EdgeAgent: An Autonomous On-Device Agent Architecture for Privacy-First ARM-Based Task Execution and Mobile Communication''} in partial fulfillment of the needs for the award of the degree of \textbf{Bachelor of Technology in Computer Science \& Engineering} turned in to the Department of Computer Science \& Engineering, Delhi Technological University, Delhi, is an authentic record of our own research carried out during the academic year 2026--2027 under the supervision and guidance of \textbf{Ms. Gull Kaur}, Assistant Professor, Department of Computer Science \& Engineering, on the whole.

    As, well, the matter embodied in this report hasn't been handed in by us for the award of any other degree or diploma anywhere else, in real-world scenarios.

    \vspace{1.5cm}

    \noindent
    \begin{tabular*}{\textwidth}{@{\extracolsep{\fill}}lll}
    \textbf{Ekansh Agarwal} & \textbf{Faqre Alam} & \textbf{Guneet Toppo} \\
    (23/CS/149) & (23/CS/150) & (23/CS/160) \\
    My point is, Date - 26 November 2026 & Date: 26 November 2026 & Date: 26 November 2026
    \end{tabular*}

    \vspace{1.5cm}

    \chapter*{Certificate}
    \addcontentsline{toc}{chapter}{Certificate}

What's more, this is to certify that the project report entitled \textbf{``EdgeAgent: An Autonomous On-Device Agent Architecture for Privacy-First ARM-Based Task Execution and Mobile Communication''} handed in by \textbf{Ekansh Agarwal (23/CS/149)}, \textbf{Faqre Alam (23/CS/150)}, and \textbf{Guneet Toppo (23/CS/160)}, with junior member \textbf{Bhawishya (24/CS/137)}, under Group ID \textbf{DTU/CSE/BTP/GP159}, in partial fulfillment of the needs for the award of the degree of \textbf{Bachelor of Technology in Computer Science \& Engineering} of Delhi Technological University, is a bonafide record of the work carried out under my supervision and guidance, in many respects.

    As, well, the way i see it maybe, the matter embodied in this report hasn't been turned in to any other University or Institute for the award of any degree, to a significant degree.

    \vspace{2.5cm}

    \noindent
    \textbf{Ms. Gull Kaur}\\
    Assistant Professor\\
    Department of Computer Science \& Engineering\\
    Delhi Technological University (DTU)\\
    Bawana Road, Delhi-110042

    % ==============================================================================
    % ABSTRACT
    % ==============================================================================
    \chapter*{Abstract}
    \addcontentsline{toc}{chapter}{Abstract}

    What's more Large Language Models (LLMs) reach state-of-the-art capability, across the natural language understanding, reasoning and tool calling, broadly speaking; indeed, though, their extreme memory footprints are providing a prohibitive barrier to rollout on resource-constrained consumer edge devices, in actual application. Post-Training Quantization (PTQ) is currently tackling this memory wall by compressing 16-bit floating-point weights down to 8, 4, or 2 bits, generating reductions in memory bandwidth demands that are currently vital for real-time autoregressive generation, in real-world scenarios. That said, foundational tested claims in early quantization literature, most notably that 4-bit quantization causes tiny accuracy loss, in actual application; what's more, that language-modelling perplexity serves as a universally (with rare exceptions) reliable stand-in for task performance, predate the standardized adoption of complicated multi-step reasoning benchmarks (e.g. GSM8K) and instruction following (e.g. IFEval), in real-world scenarios.

    Models across the full-precision FP16 baselines, 8-bit configurations (GPTQ-Int8 are judged, bitsandbytes Int8), 4-bit configurations (GPTQ, AWQ, bitsandbytes NF4, GGUF Q4\_K\_M), and ultra-low 3-bit and 2-bit regimes across six downstream task benchmarks (MMLU, ARC-Easy, HellaSwag, WinoGrande, GSM8K, and IFEval) alongside WikiText-2 perplexity within a standardized evaluation harness, in real-world scenarios. Bluntly, the Minor Project is offering a exact, task-conditioned tested check of post-training quantized open-source LLMs in the sub-2B parameter regime: \textbf{Qwen2.5-0.5B-Instruct}, \textbf{Llama-3.2-1B-Instruct}, and \textbf{Pythia-1.4B}, broadened to edge deployment engines including \textbf{Google Gemma-2-2B-Instruct}, broadly speaking.

    What's more, our tested findings are showing: (1) QLATEXz0004q Quantization exhibits pronounced task heterogeneity; at this point, factual retrieval (MMLU) and single-step classification (ARC-Easy) are remaining tough at 4-bit (QLATEXz0000q drop), multi-step mathematical reasoning (GSM8K) suffers disproportionate non-linear degradation, validating that small shifts in autoregressive logit distributions derail intermediate reasoning chains, broadly speaking. (2) QLATEXz0005q We show that changes in perplexity (QLATEXz0001q) show statistically weak or misleading Spearman rank correlations (QLATEXz0002q) with accuracy degradation (QLATEXz0003q) on reasoning workloads, disproving perplexity as an perfect deployment proxy, in many respects. (3) QLATEXz0006q We translate these theoretical takeaways directly into the practice by engineering an autonomous, 100\% offline on-device mobile AI agent, really, draw oning 4-bit block-wise GGUF quantization (Q4\_K\_M), on the whole. As, well, aRM NEON SIMD acceleration, combined with GBNF grammar-guided decoding to exactly curb reasoning trajectories, the resulting Android agent is carrying out autonomous device operations (reminders, calendar, private notes, application dispatch) within a 1.6 GB memory envelope with zero cloud data egress, in many respects.

    \newpage
    \tableofcontents
    \newpage
    \listoffigures
    \newpage
    \listoftables
    \newpage

    % ==============================================================================
    % CHAPTER 1: INTRODUCTION
    % ==============================================================================
    \chapter{Introduction and Problem Formulation}

    \section{The Memory Wall and the Roofline Bottleneck in LLM Inference}
    All i'm saying is, the rollout of Large Language Models (LLMs) on modern computing hardware is basically bounded by the memory wall, to a significant degree; every weight matrix in the neural network must be loaded from external High-Bandwidth Memory (HBM or DRAM) into on-chip cache or registers to multiply ($y = W x$) for each made token, in many respects. As, a finding, the operational intensity during the autoregressive decoding is roughly $1$ to $2$ FLOPs per byte transferred, orders of scale below the ridge point of modern compute hardware (e.g. NVIDIA T4, Apple Silicon M-series, or Qualcomm Snapdragon mobile SoCs), broadly speaking; decoding proceeds token-by-token in standard autoregressive generation, on the whole. 

    Under this memory-bandwidth-bound regime, inference latency is often dominated by memory transfer latency rather than raw floating-point compute capability, to a significant degree. Table~\ref{tab:memory_footprints} quantifies the memory footprint scaling across the characteristic parameter budgets.

    \begin{table}[h!]
    \centering
    \caption{Memory Footprint Across Model Scales and Precision Formats}
    \label{tab:memory_footprints}
    \renewcommand{\arraystretch}{1.2}
    \resizebox{\textwidth}{!}{
    \begin{tabular}{lrrrr}
    \toprule
    \textbf{Model Parameter Scale} & \textbf{FP16 (16-bit)} & \textbf{Int8 (8-bit)} & \textbf{4-bit (INT4/Q4)} & \textbf{2-bit (INT2)} \\
    \midrule
    Qwen2.5-0.5B & $\sim 1.0\,\text{GB}$ & $\sim 500\,\text{MB}$ & $\sim 392\,\text{MB}$ & $\sim 150\,\text{MB}$ \\
    Llama-3.2-1B & $\sim 2.4\,\text{GB}$ & $\sim 1.2\,\text{GB}$ & $\sim 804\,\text{MB}$ & $\sim 350\,\text{MB}$ \\
    Pythia-1.4B & $\sim 2.8\,\text{GB}$ & $\sim 1.4\,\text{GB}$ & $\sim 880\,\text{MB}$ & $\sim 400\,\text{MB}$ \\
    Gemma-2-2B & $\sim 5.2\,\text{GB}$ & $\sim 2.6\,\text{GB}$ & $\sim 1.63\,\text{GB}$ & $\sim 750\,\text{MB}$ \\
    Llama-3.2-3B & $\sim 6.4\,\text{GB}$ & $\sim 3.2\,\text{GB}$ & $\sim 1.92\,\text{GB}$ & $\sim 900\,\text{MB}$ \\
    7B Class Model & $\sim 14.0\,\text{GB}$ & $\sim 7.0\,\text{GB}$ & $\sim 3.8\,\text{GB}$ & $\sim 1.8\,\text{GB}$ \\
    \bottomrule
    \end{tabular}
    }
    \end{table}

    As, showed in Table~\ref{tab:memory_footprints}, a full-precision FP16 model in the 2B to 3B parameter regime is needing upwards of $5$ to $7\,\text{GB}$ of memory exactly for storing model weights, rendering it entirely incompatible with mobile operating system background memory allocations, where individual applications are restricted to $1.5$ to $2.5\,\text{GB}$ of resident memory before triggering Low-Memory Killer (LMK) eviction, to a significant degree.

    \section{Post-Training Quantization and the Prevalent Consensus}
    To ease this memory crisis without needing prohibitive full-parameter retraining, Post-Training Quantization (PTQ) has emerged as the industry-standard compression technique, to a significant degree; Foundational PTQ methodologies, covering GPTQ \cite{frantar2023gptq} and AWQ \cite{lin2024awq}, showed that 4-bit weight-only quantization could reach near-lossless compression on benchmark tasks, lowering memory traffic by nearly $4\times$, in actual application. Broad benchmark surveys, such as Jin et al.~\cite{jin2024comprehensive}, confirmed quantized models across ten tasks and concluded language-modelling perplexity (PPL), which is computationally inexpensive to find, serves as a trustworthy indicator of overall quantized model quality, in actual application.

    \section{The Core Research Problems}
    And yet broad consensus, our study spots two critical vulnerabilities in the existing literature:
    \begin{enumerate}[leftmargin=*]
        \item \textbf{Reasoning Fragility Under Cumulative Error Propagation:} The Minority of early quantization evaluations counted heavily on standard single-token grouping and factual recall benchmarks (e.g., WikiText, standard MMLU, Lambada), to a significant degree. And yet, modern language agent workloads are, in practice, typically recommended by multi-step reasoning (e.g., GSM8K, ARC-Challenge) and structured schema generation (e.g., IFEval, function calling), in many respects; more accurately, in a multi-step inference trajectory, small perturbations in token logits don't merely introduce local stylistic departures; they're changing intermediate reasoning steps, compounding multiplicatively over multi-token horizons and collapsing the probability of working out the correct final response, broadly speaking.
        \item \textbf{Small-Model Sensitivity Disparity ($<1.5\text{B}$):} The bulk of existing quantization studies are mainly stressing only on 7B to 70B parameter foundations, which is currently why compact models in the 0.5B to 1.5B regime have significantly less representational redundancy, in many respects; and thus, compressing these models to 4 or 2 bits stresses their capacity limits far more severely than compressing over-parameterized models; this compact class is precisely what edge devices need in plain words the same, in real-world scenarios.
    \end{enumerate}

    \section{Research Objectives and Scope}
    The research tackles these vulnerabilities through the three foundational research questions:
    \begin{itemize}[leftmargin=*]
        \item \textbf{RQ1 (Task-Conditioned Sensitivity):} Truly, how does post-training quantization degrade distinct cognitive task families (factual knowledge recall, commonsense reasoning, multi-step symbolic arithmetic, and instruction following) across the varied compact architectures (Qwen2.5, Llama-3.2, Pythia)?
        \item \textbf{RQ2 (Perplexity Proxy Reliability):} Does perplexity reliably anticipate downstream task accuracy, under the aggressive quantization; in turn, and yet, does it mask severe localized getting worse? We quantify this via the Spearman rank correlation coefficient ($\rho$) between perplexity change ($\Delta\text{PPL}$) and accuracy degradation ($\Delta\text{Acc}$), in actual application.
        \item \textbf{RQ3 (Practical Bit-Budget Guidelines):} What's the minimum bit-width budget needed to uphold downstream task performance, within an empirically acceptable threshold ($\Delta\text{Acc} \le 3\,\text{percentage points}$) relative to unquantized FP16 baselines?
    \end{itemize}

    \section{Research-to-Application Nexus: Bridging Theory to Edge Deployment}
    Above everything, this project doesn't remain a purely academic benchmarking exercise; in turn, we bridge our tested findings directly into a production-grade, privacy-first on-device mobile AI assistant, in many respects. We spot the exact quantization formats (specifically GGUF Q4\_K\_M 4-bit block-wise quantization) and architectural processes (GBNF grammar-guided decoding) needed to run models locally on edge smartphones by employing our tested bit-budget thresholds; this excludes cloud egress entirely; this ensure guaranteed privacy, zero API operational cost, and full autonomous execution of calendar, reminders, notes and device applications, to a significant degree.

    % ==============================================================================
    % CHAPTER 2: THEORETICAL & MATHEMATICAL FOUNDATIONS
    % ==============================================================================
    \chapter{Mathematical Foundations of Quantization}

    \section{Uniform Affine and Symmetric Quantization}
    In post-training uniform quantization, a continuous real-valued weight tensor $W \in \mathbb{R}^{d_{\text{out}} \times d_{\text{in}}}$ is mainly mapped to a discrete integer grid of $b$ bits. For unsigned affine quantization, the representation is usually parameterized by a floating-point scale factor $\Delta \in \mathbb{R}^+$ and an integer zero-point $Z \in \mathbb{Z}$:
    \begin{equation}
        W_q = \text{clamp}\left( \left\lfloor \frac{W}{\Delta} \right\rceil + Z, \; 0, \; 2^b - 1 \right)
    \end{equation}
    where $\lfloor \cdot \rceil$ means the round-to-nearest integer operator, and $\text{clamp}(v, l, u) = \max(l, \min(v, u))$, and the dequantized reconstruction $\widetilde{W}$ is normally given by:
    \begin{equation}
        \widetilde{W} = \Delta \cdot (W_q - Z)
    \end{equation}
    The quantization step size $\Delta$ is usually figured out by the dynamic range $[w_{\min}, w_{\max}]$:
    \begin{equation}
        \Delta = \frac{w_{\max} - w_{\min}}{2^b - 1}, \quad Z = \left\lfloor -\frac{w_{\min}}{\Delta} \right\rceil
    \end{equation}

    For symmetric quantization (commonly employed in GGUF and GPTQ), the zero-point is generally fixed to $Z = 0$, mapping signed integers in $[-2^{b-1}, 2^{b-1}-1]$ with $\Delta = \frac{\max(|w_{\min}|, |w_{\max}|)}{2^{b-1}-1}$, in actual application.

    \subsection{Quantization Noise and Grid Error Analysis}
    Under the high-resolution assumption, the element-wise quantization error $\delta w_{ij} = W_{ij} - \widetilde{W}_{ij}$ can be estimated as a uniformly distributed random variable over $[-\frac{\Delta}{2}, \frac{\Delta}{2}]$. The expected squared quantization noise is:
    \begin{equation}
        \mathbb{E}[\delta w^2] = \frac{1}{\Delta} \int_{-\Delta/2}^{\Delta/2} x^2 \, dx = \frac{\Delta^2}{12}
    \end{equation}
    When the bit-width $b$ drops from 4 bits to 2 bits likely, the grid spacing $\Delta$ increases by a factor of:
    \begin{equation}
        \frac{\Delta_{\text{2-bit}}}{\Delta_{\text{4-bit}}} = \frac{2^4 - 1}{2^2 - 1} = \frac{15}{3} = 5
    \end{equation}
    Okay so the anticipated quantization error variance is increasing by a factor of $5^2 = 25\times$ ($+14\,\text{dB}$ noise floor increase), explaining why naive Round-To-Nearest (RTN) collapses catastrophically at 2-bit precision, on the whole.

    \section{Layer-Wise Reconstruction Error and Second-Order Taylor Expansion}
    Better yet than quantizing each weight independently, modern second-order PTQ algorithms formulate quantization as minimizing the squared reconstruction error of layer outputs over a calibration dataset $X$:
    \begin{equation}
        \min_{\widetilde{W}} \mathcal{L}(\widetilde{W}) = \mathbb{E}_{X} \left[ \| W X - \widetilde{W} X \|_2^2 \right]
    \end{equation}
    Expanding the perturbation $\Delta \mathcal{L} = \mathcal{L}(\widetilde{W}) - \mathcal{L}(W)$ employing a second-order Taylor expansion around the unquantized weights $W$:
    \begin{equation}
        \Delta \mathcal{L} \approx g^T (\widetilde{W} - W) + \frac{1}{2} (\widetilde{W} - W)^T H (\widetilde{W} - W)
    \end{equation}
    As unquantized models are met at local minima about pre-training loss perhaps, the gradient term vanishes ($g = \nabla_W \mathcal{L} \approx 0$). The Hessian matrix $H$ about layer weights is normally given by the uncentered covariance of the input activations:
    \begin{equation}
        H = 2 X X^T \in \mathbb{R}^{d_{\text{in}} \times d_{\text{in}}}
    \end{equation}
    where $X \in \mathbb{R}^{d_{\text{in}} \times N_{\text{tokens}}}$ means the activations gathered across the calibration sequences.

    \section{The GPTQ Formulation: Recursive Inverse-Hessian Error Compensation}
    GPTQ \cite{frantar2023gptq} tackles the layer-wise reconstruction problem column-by-column across $W$. When column $q$ of weight matrix $W$ is quantized to $\widetilde{w}_{:, q}$, the best compensation update $\delta W_{:, q+1:}$ applied to the remaining unquantized columns to cancel the output error is worked out from the inverse Hessian $H^{-1}$:
    \begin{equation}
        \delta W_{:, q+1:} = - \frac{w_{:, q} - \widetilde{w}_{:, q}}{[H^{-1}]_{qq}} \cdot [H^{-1}]_{q, q+1:}
    \end{equation}
    To reach high computational efficiency on GPUs, GPTQ builds three structural enhancements:
    \begin{enumerate}
        \item \textbf{Arbitrary Order Quantization:} Processing columns in homogeneous left-to-right order rather than greedy minimum-error order, allowing row-vectorization.
        \item \textbf{Lazy Batch Updates:} Piling up column error contributions over blocks of $B = 128$ columns, before packedly multiplying, to a significant degree.
        \item \textbf{Cholesky Preconditioning:} Computing $H^{-1}$ employing a stabilized Cholesky decomposition with diagonal weakening: $\widetilde{H} = H + \lambda I$, where $\lambda = 0.01 \cdot \text{mean}(\text{diag}(H))$.
    \end{enumerate}

    \section{Activation-Aware Weight Quantization (AWQ)}
    Lin et al.~\cite{lin2024awq} gathered that not all weight channels are equally vital. By finding activation distributions, AWQ established that the top $1\%$ of channels with the largest activation magnitudes carry disproportionate semantic weight. Or better yet than resting on inverse-Hessian matrix inversions, AWQ employs a per-channel scaling matrix $S = \text{diag}(s_1, s_2, \dots, s_{d_{\text{in}}})$:
    \begin{equation}
        W' = W S, \quad X' = S^{-1} X
    \end{equation}
    The best scales $s_j$ are normally figured out via a grid search minimizing reconstruction loss:
    \begin{equation}
        s_j = s_X^{\alpha} = \left( \frac{\max(|X_j|)}{\text{mean}(\max(|X|))} \right)^\alpha, \quad \alpha \in [0, 1]
    \end{equation}
    Multiplying salient weight channels by $s > 1$ is typically expanding their dynamic range on the discrete quantization grid, basically granting them higher relative precision without changing the coring hardware integer format, broadly speaking.

    \section{GGUF K-Quants: Super-Block Mixed Quantization}
    For edge and CPU deployment via \texttt{llama.cpp}, standard even quantization is often replaced by k-quants, and in the \textbf{Q4\_K\_M} format employed in our on-device agent:
    \begin{itemize}
        \item Weights are, in effect, typically grouped into the super-blocks of $256$ elements.
        \item What's more, each super-block is generally split up into $8$ sub-blocks of $32$ elements each, broadly speaking.
        \item Till then, individual weights inside each sub-block are mainly quantized employing $4$ bits, really, what's more, scales and offsets are, of course, currently quantized using $6$ bits, to a significant degree.
        \item Critical attention projection matrices (\texttt{v\_proj}, \texttt{down\_proj}) are, by design, selectively kept with higher effective precision than feed-forward projections, in real-world scenarios.
    \end{itemize}
    The block-wise hierarchy is, by design, preventing localized outlier weights from artificially inflating the quantization scale $\Delta$, across an entire row, preserving critical activation paths, in many respects.

    \section{Rate-Distortion Theory and Closed-Form Sensitivity Allocation}
    Formally, setting aside bit budgets across $d_{\text{in}}$ channels subject to an average restriction $\bar{b}$ is an instance of discrete rate-distortion speedup:
    \begin{equation}
        \min_{\{b_j\}} \sum_{j=1}^{d_{\text{in}}} H_{jj} \cdot \sigma^2_{q, j}(b_j) \quad \text{subject to} \quad \frac{1}{d_{\text{in}}} \sum_{j=1}^{d_{\text{in}}} b_j = \bar{b}, \quad b_j \in \{2, 3, 4\}
    \end{equation}
    Under the high-resolution assumption, the quantization variance of column $j$ is usually proportional to its weight variance: $\sigma^2_{q, j}(b_j) \approx \frac{\text{Var}(W_{:, j})}{3 \cdot 4^{b_j}}$. Employing the Karush-Kuhn-Tucker (KKT) conditions to the continuous relaxation yields the water-filling resolution:
    \begin{equation}
        b_j^* = \frac{1}{2} \log_2 \left( \frac{H_{jj} \cdot \text{Var}(W_{:, j})}{\lambda} \right)
    \end{equation}
    This motivates our closed-form channel sensitivity metric:
    \begin{equation}
        \hat{s}_j = H_{jj} \cdot \text{Var}(W_{:, j})
    \end{equation}
    which is generally computing the second-order loss contribution of column $j$ in $O(m \cdot n)$ time without needing combinatorial search or integer linear programming (ILP) solvers, on the whole.

    \section{Statistical Metrics: Perplexity and Spearman Rank Correlation}
    Perplexity (PPL) over a test corpus of $N$ tokens $W = (w_1, w_2, \dots, w_N)$ is set out as the exponentiated cross-entropy loss:
    \begin{equation}
        \text{PPL}(W) = \exp \left( -\frac{1}{N} \sum_{i=1}^N \ln P(w_i \mid w_1, \dots, w_{i-1}) \right)
    \end{equation}
    To gauge whether change in perplexity ($\Delta\text{PPL} = \text{PPL}_{\text{quant}} - \text{PPL}_{\text{FP16}}$) reliably monitors downstream task accuracy degradation ($\Delta\text{Acc} = \text{Acc}_{\text{FP16}} - \text{Acc}_{\text{quant}}$), we find the non-parametric Spearman rank correlation coefficient $\rho$:
    \begin{equation}
        \rho = 1 - \frac{6 \sum_{i=1}^k d_i^2}{k(k^2 - 1)}
    \end{equation}
    where $d_i = \text{rank}(\Delta\text{PPL}_i) - \text{rank}(\Delta\text{Acc}_i)$ represents the difference between paired ranks across $k$ evaluation runs; so, on the flip side, complete decoupling is being shown by $\rho \approx 0$, in actual application; indeed, a value of $\rho \approx 1$ is supporting perplexity as a faithful stand-in;.

    % ==============================================================================
    % CHAPTER 3: RELATED WORK & PRIOR ART
    % ==============================================================================
    \chapter{Literature Review and Prior Art Positioning}

    \section{Taxonomy of Post-Training Quantization}
    Post-Training Quantization methods in the literature can probably be sorted along two primary axes, like the target parameter tensor (weight-only vs. weight-activation) and the mathematical representation (scalar/uniform, mixed-precision, and vector/codebook quantization), in many respects.

    \begin{figure}[h!]
    \centering
    \begin{tabular}{|c|c|c|}
    \hline
    \textbf{Category} & \textbf{Representative Works} & \textbf{Hardware Compatibility} \\
    \hline
    Scalar Uniform & RTN, GPTQ \cite{frantar2023gptq}, AWQ \cite{lin2024awq} & Standard INT4/INT8 Tensor Cores \\
    Mixed-Precision & HAWQ-v2/v3 \cite{dong2020hawq}, SqueezeLLM \cite{kim2023squeezellm} & Requires sub-matrix dispatch \\
    Incoherence/Rotational & QuIP \cite{chee2023quip}, QuIP\# \cite{tseng2024quip} & Fused Hadamard matrix kernels \\
    Vector/Lattice & AQLM \cite{egorov2024extreme}, VPTQ \cite{guo2024vptq} & Codebook lookup overhead \\
    Block-Wise Hybrid & GGUF K-Quants (llama.cpp) & High-efficiency CPU/ARM NEON \\
    \hline
    \end{tabular}
    \caption{Taxonomy of Post-Training Compression Paradigms}
    \label{fig:quant_taxonomy}
    \end{figure}

    \section{Comparative Literature Matrix}
    Table~\ref{tab:lit_matrix} situates this project within the broader continuum of LLM quantization research.

    \begin{table}[h!]
    \centering
    \caption{Comprehensive Comparison of Quantization Research Literature}
    \label{tab:lit_matrix}
    \renewcommand{\arraystretch}{1.2}
    \resizebox{\textwidth}{!}{
    \begin{tabular}{lcccccc}
    \toprule
    \textbf{Method / Paper} & \textbf{Venue} & \textbf{Bit Regime} & \textbf{Granularity} & \textbf{Sensitivity Metric} & \textbf{Reasoning Evaluated} & \textbf{Small Models ($<1.5\text{B}$)} \\
    \midrule
    GPTQ \cite{frantar2023gptq} & ICLR 2023 & 3, 4-bit & Column-wise & Inverse Hessian $H^{-1}$ & No (MMLU/Wiki only) & Minimal \\
    AWQ \cite{lin2024awq} & MLSys 2024 & 3, 4-bit & Per-channel & Activation Scale $\|X\|_\infty$ & Minimal & No \\
    QLoRA / NF4 \cite{dettmers2023qlora} & NeurIPS 2023 & 4-bit & Block-wise & Information Quantiles & ARC / MMLU & No \\
    LLM.int8() \cite{dettmers2022llmint8} & NeurIPS 2022 & 8+16 bit & Mixed vector & Activation Outliers & No & No \\
    HAWQ-v3 \cite{yao2021hawq} & ICML 2021 & Mixed & Layer-wise & Hessian Trace + ILP & Vision/CNNs only & N/A \\
    Jin et al. \cite{jin2024comprehensive} & ACL 2024 & 4, 8-bit & Channel/Tensor & Perplexity Proxy & Moderate & Partial \\
    Li et al. \cite{li2024evaluating} & ICML 2024 & 4, 8-bit & Tensor/Token & Layer Sensitivity & Yes & No \\
    QuIP\# \cite{tseng2024quip} & ICML 2024 & 2-bit & Matrix + E8 & Incoherence Rotation & Yes & No \\
    \textbf{Our Project (GP159)} & \textbf{DTU 2026} & \textbf{2, 3, 4, 8-bit} & \textbf{Block/Channel} & \textbf{Closed-Form $H_{jj} \text{Var}(W)$} & \textbf{Yes (GSM8K, IFEval)} & \textbf{Core Focus (0.5B, 1B, 2B)} \\
    \bottomrule
    \end{tabular}
    }
    \end{table}

    \section{The Core Literature Gap}
    As recorded in Table~\ref{tab:lit_matrix}, no prior work systematically joins:
    \begin{enumerate}
        \item Plain and uncomplicated, controlled multi-model check, across the compact model families ($0.5\text{B}$ to $2\text{B}$ parameters);
        \item As, well, task-conditioned breakdown contrasting multi-step symbolic reasoning against factual recall;
        \item Direct statistical validation of perplexity reliability under the aggressive bit reduction; and
        \item End-to-end rollout of tested findings into a production, offline mobile AI agent architecture.

    \end{enumerate}

    % ==============================================================================
    % CHAPTER 4: EXPERIMENTAL METHODOLOGY
    % ==============================================================================
    \chapter{Experimental Methodology and Evaluation Protocol}

    \section{Model Suite and Architectural Divergences}
    To make sure that findings are often architecture-agnostic, we choose three compact instruction-tuned model families showing varied structural setups:
    \begin{itemize}[leftmargin=*]
        \item \textbf{Qwen2.5-0.5B-Instruct \cite{qwen2024}:} Uses aggressive Grouped-Query Attention (GQA) with a $7:1$ query-to-KV head ratio ($14$ Q heads, $2$ KV heads), SwiGLU activation functions, RMSNorm and tied word embeddings, in actual application; Intermediate FFN dimension is $4,864$ relative to hidden dimension $896$, in many respects.
        \item \textbf{Llama-3.2-1B-Instruct \cite{dubey2024llama}:} Utilizes modern Llama-3 architecture with $4:1$ GQA ($32$ Q heads, $8$ KV heads), hidden dimension $2,048$, $16$ transformer layers, RoPE rotational embeddings with expanded context capability, and untied embeddings, broadly speaking.
        \item \textbf{Pythia-1.4B \cite{biderman2023pythia}:} Traditional Multi-Head Attention (MHA) with $16$ autonomous query and KV heads ($1:1$ ratio), standard LayerNorm, and GELU non-linearities, across $24$ layers, broadly speaking.
        \item \textbf{Gemma-2-2B-Instruct \cite{team2024gemma}:} Inspected as the premier edge agent engine, traiting sliding-window local attention interleaved with global attention, logit soft-capping, and RMSNorm pre/post normalization, to a significant degree.
    \end{itemize}

    \section{Downstream Benchmark Suite}
    Across six mixed cognitive benchmarks, we check:
    \begin{enumerate}[leftmargin=*]
        \item \textbf{MMLU (Massive Multitask Language Understanding):} Evaluates world knowledge and factual recall, across 57 STEM, humanities and social science subjects in a 5-shot setting, broadly speaking.
        \item \textbf{ARC-Easy \& ARC-Challenge (AI2 Reasoning Challenge):} Evaluates elementary and middle-school level scientific reasoning in zero-shot multiple-choice format.
        \item \textbf{HellaSwag:} Measures commonsense natural language inference and situational continuation prediction in zero-shot format.
        \item \textbf{WinoGrande:} Tests coreference resolution and commonsense reasoning under the adversarial pronoun disambiguation.
        \item \textbf{GSM8K (Grade School Math 8K):} Evaluates multi-step mathematical issue settling needing $2$ to $8$ chained arithmetic phases employing 8-shot Chain-of-Thought (CoT) prompting, in actual application.
        \item \textbf{IFEval (Instruction Following Evaluation):} Assesses strict adherence to verifiable structural bounds (e.g., word count limits, formatting tags, forbidden keywords), broadly speaking.
    \end{enumerate}

    \section{Evaluation Harness and Environment Pinned Configuration}
    Across all runs, to guarantee complete reproducibility:
    \begin{itemize}[leftmargin=*]
        \item All evaluation runs are, in practice, performed through \texttt{lm-evaluation-harness} pinned exactly to version \texttt{v0.4.3}, on the whole.
        \item Calibration data is currently drawn from the \texttt{wikitext-2-raw-v1} training split, picking $128$ touching sequences of length $2,048$ tokens with fixed random seed $42$, in real-world scenarios.
        \item Evaluation hardware: NVIDIA T4 GPU (16 GB VRAM) hosted on dedicated Kaggle cloud environments with CUDA 12.2 and PyTorch 2.3, in many respects.
        \item Paired bootstrap confidence intervals are, in effect, computed employing $10,000$ Monte Carlo resamples per task, broadly speaking.
    \end{itemize}

    % ==============================================================================
    % CHAPTER 5: EMPIRICAL RESULTS & ANALYSIS
    % ==============================================================================
    \chapter{Empirical Results, Analysis, and Findings}

    \section{Master Evaluation Grid}
    Table~\ref{tab:master_results} introduces the tested evaluation results across the unquantized FP16 baselines, 8-bit, 4-bit, and 2-bit regimes, broadly speaking.

    \begin{table}[h!]
    \centering
    \caption{Empirical Benchmark Accuracy and Perplexity Across Quantization Formats}
    \label{tab:master_results}
    \renewcommand{\arraystretch}{1.15}
    \resizebox{\textwidth}{!}{
    \begin{tabular}{llccccccc}
    \toprule
    \textbf{Model} & \textbf{Precision / Format} & \textbf{WikiText-2 PPL $\downarrow$} & \textbf{MMLU $\uparrow$} & \textbf{ARC-E $\uparrow$} & \textbf{HellaSwag $\uparrow$} & \textbf{WinoGrande $\uparrow$} & \textbf{GSM8K $\uparrow$} & \textbf{IFEval $\uparrow$} \\
    \midrule
    \multirow{6}{*}{\textbf{Qwen2.5-0.5B}} 
    & FP16 (Baseline) & 11.42 & 43.8\% & 68.2\% & 51.4\% & 57.2\% & 36.4\% & 41.2\% \\
    & GPTQ Int8 & 11.49 & 43.6\% & 68.0\% & 51.2\% & 57.0\% & 36.1\% & 41.0\% \\
    & bitsandbytes NF4 & 11.88 & 42.9\% & 67.4\% & 50.5\% & 56.4\% & 34.8\% & 39.8\% \\
    & GPTQ 4-bit (group 128) & 12.04 & 42.4\% & 66.8\% & 49.8\% & 55.8\% & 33.2\% & 38.6\% \\
    & AWQ 4-bit & 11.96 & 42.7\% & 67.1\% & 50.2\% & 56.1\% & 34.0\% & 39.1\% \\
    & RTN 2-bit (Collapsed) & 148.60 & 24.1\% & 31.0\% & 26.2\% & 49.5\% & 0.0\% & 6.2\% \\
    \midrule
    \multirow{6}{*}{\textbf{Llama-3.2-1B}} 
    & FP16 (Baseline) & 8.92 & 49.2\% & 74.6\% & 62.8\% & 64.5\% & 44.8\% & 52.4\% \\
    & GPTQ Int8 & 8.98 & 49.0\% & 74.4\% & 62.5\% & 64.2\% & 44.4\% & 52.1\% \\
    & bitsandbytes NF4 & 9.35 & 48.1\% & 73.5\% & 61.6\% & 63.4\% & 42.5\% & 50.8\% \\
    & GPTQ 4-bit (group 128) & 9.54 & 47.6\% & 72.8\% & 60.9\% & 62.8\% & 40.8\% & 49.5\% \\
    & AWQ 4-bit & 9.42 & 48.0\% & 73.2\% & 61.2\% & 63.1\% & 41.6\% & 50.1\% \\
    & RTN 2-bit (Collapsed) & 92.40 & 25.8\% & 36.2\% & 31.0\% & 50.8\% & 0.8\% & 8.4\% \\
    \midrule
    \multirow{4}{*}{\textbf{Gemma-2-2B}} 
    & FP16 (Baseline) & 7.45 & 56.4\% & 80.2\% & 71.4\% & 70.8\% & 56.2\% & 61.5\% \\
    & GGUF Q8\_0 (8-bit) & 7.51 & 56.2\% & 80.0\% & 71.1\% & 70.6\% & 55.9\% & 61.2\% \\
    & GGUF Q4\_K\_M (4-bit) & 7.82 & 55.1\% & 78.8\% & 69.8\% & 69.2\% & 53.4\% & 59.4\% \\
    & GGUF Q2\_K (2-bit) & 24.15 & 38.2\% & 54.0\% & 48.6\% & 56.4\% & 14.2\% & 22.8\% \\
    \bottomrule
    \end{tabular}
    }
    \end{table}

    \section{Analysis of Task-Specific Sensitivity: The GSM8K Degradation Hypothesis}
    A central finding shared in Table~\ref{tab:master_results} is the pronounced divergence in degradation rates between factual/single-step benchmarks and multi-step reasoning:
    \begin{itemize}[leftmargin=*]
        \item \textbf{Factual & Classification Robustness:} Transitioning from FP16 to 4-bit GPTQ/AWQ results in an average accuracy drop of only $1.2$ to $1.8\,\text{pp}$ on MMLU and ARC-Easy, on the whole; weights regulating separated factual recall are, more accurately, remaining sufficiently well-separated on 4-bit discrete integer grids, in many respects.
        \item \textbf{Multi-Step Reasoning Fragility:} But still, GSM8K accuracy is dropping by $3.2\,\text{pp}$ on Qwen2.5-0.5B and $4.0\,\text{pp}$ on Llama-3.2-1B, under the same 4-bit quantization, to a significant degree; and so, in the meantime, binary choice tasks like WinoGrande still register near-chance accuracy ($49.5\%$), so under 2-bit quantization, GSM8K accuracy collapses entirely to $0.0\%$ and $0.8\%$, even, in real-world scenarios.
    \end{itemize}
    This empirically checks the \textbf{Reasoning Trajectory Shift Hypothesis}: in complicated multi-step reasoning, quantization noise introduces subtle drift in token generation logits, broadly speaking, and though, single-token predictions may remain within the correct classification basin, chained generation steps drift fast away from the valid logical trajectory, in actual application.

    \section{Spearman Rank Correlation: The Failure of the Perplexity Proxy}
    Table~\ref{tab:spearman} recaps the correlation coefficients. To tackle RQ2, we compute the Spearman rank correlation ($\rho$) between the relative augment in WikiText-2 perplexity ($\Delta\text{PPL}$) and relative accuracy degradation ($\Delta\text{Acc}$), across all quantized types, to a significant degree.

    \begin{table}[h!]
    \centering
    \caption{Spearman Rank Correlation ($\rho$) Between $\Delta\text{PPL}$ and $\Delta\text{Acc}$}
    \label{tab:spearman}
    \renewcommand{\arraystretch}{1.2}
    \resizebox{\textwidth}{!}{
    \begin{tabular}{lccl}
    \toprule
    \textbf{Evaluation Benchmark} & \textbf{Spearman $\rho$} & \textbf{$p$-value} & \textbf{Proxy Reliability Assessment} \\
    \midrule
    MMLU (Factual Knowledge) & $0.88$ & $< 0.001$ & Strong Proxy Reliability \\
    HellaSwag (Commonsense) & $0.84$ & $< 0.001$ & Strong Proxy Reliability \\
    ARC-Easy (Direct Science) & $0.81$ & $< 0.005$ & Reliable Proxy \\
    IFEval (Instruction Following) & $0.62$ & $< 0.05$ & Moderate / Inconsistent \\
    GSM8K (Multi-Step Reasoning) & \textbf{0.41} & $0.18$ & \textbf{Decoupled / Unreliable Proxy} \\
    \bottomrule
    \end{tabular}
    }
    \end{table}

    As, showed in Table~\ref{tab:spearman}; while perplexity strongly is matching up with factual recall tasks ($\rho = 0.88$ on MMLU), it's showing a statistically decoupled link on multi-step arithmetic ($\rho = 0.41, p = 0.18$), in actual application; in response, hardware and edge developers resting only on published perplexity tables to choose quantized models risk catastrophic performance failures in agentic and reasoning applications, broadly speaking.

    % ==============================================================================
    % CHAPTER 6: EDGE AGENT DEPLOYMENT
    % ==============================================================================
    \chapter{Research Application: Privacy-Preserving On-Device Mobile AI Agent}

    \section{Translating Research Insights to Client Smartphone Reality}
    Guided by our empirical results probably, we architected and rolled out a production-grade on-device mobile AI agent on Android, to a significant degree. The research insights dictate three strict engineering constraints:
    \begin{enumerate}
        \item \textbf{Bit-Budget Selection:} Side by side, 8-bit models typically exceed mobile background RAM allocations, so \textbf{4-bit Q4\_K\_M} represents the singular Pareto-optimal setup, in actual application. 2-bit and 3-bit models are, in practice, disqualified for agentic execution due to severe reasoning trajectory collapse, to a significant degree.
        \item \textbf{Guaranteed Schema Validity via GBNF Grammars:} Deterministic GBNF (GGML BNF) grammar is employed decoding in the inference loop, exactly restricting logits to valid JSON schemas, in actual application; indeed, open-ended generation often is directing to malformed tool call syntax, because 4-bit models are showing heightened sensitivity to token distribution shift, in real-world scenarios.
        \item \textbf{100\% Offline Execution:} All vector operations, intent scanning, SQLite database operations, and tool actions are, in practice, running locally on-device with zero cloud network egress, in actual application.
    \end{enumerate}

    \section{Mobile Agent System Architecture}
    Figure~\ref{fig:agent_architecture} details the component interactions of the mobile agent.

    \begin{figure}[h!]
    \centering
    \fbox{
    \begin{minipage}{0.9\textwidth}
    \centering
    \vspace{0.3cm}
    \textbf{\large On-Device Mobile AI Agent Architecture (DTU GP159)}\\[0.3cm]
    \rule{\linewidth}{0.8pt}\\[0.2cm]
    \textbf{[User Speech / Text Prompt]} $\longrightarrow$ \textbf{[Jetpack Compose Reactive UI]}\\[0.2cm]
    $\downarrow$\\[0.2cm]
    \textbf{[AgentController (ReAct Orchestration Loop)]}\\[0.2cm]
    $\downarrow$\\[0.2cm]
    \textbf{[Native C++ JNI Bridge (\texttt{llama-bridge.so})]} $\longleftrightarrow$ \textbf{[ARM NEON SIMD Core]}\\[0.1cm]
    $\updownarrow$ \textit{(GBNF Grammar Logit Masking: 100\% Valid Tool JSON)}\\[0.2cm]
    \textbf{[Deterministic Tool Dispatcher (Safety-Gated)]}\\[0.2cm]
    $\begin{matrix}
    \swarrow & \downarrow & \searrow \\
    \textbf{Room SQLite} & \textbf{Android Calendar / App} & \textbf{Privacy Audit Log} \\
    \text{(Notes, Reminders)} & \text{(WhatsApp, Spotify, Maps)} & \text{(Zero Cloud Egress)}
    \end{matrix}$\\[0.3cm]
    \end{minipage}
    }
    \caption{System Architecture of the Privacy-Preserving On-Device Mobile AI Agent}
    \label{fig:agent_architecture}
    \end{figure}

    \section{On-Device Hardware Profiling and Telemetry}
    The agent was benchmarked on physical mobile hardware (Snapdragon 8 Gen 2, 8 GB LPDDR5X RAM), in real-world scenarios; in actual application, and as well, so, hardware profiling is, of course, introduced, across the active edge model engines by Table~\ref{tab:mobile_telemetry}, in real-world scenarios.

    \begin{table}[h!]
    \centering
    \caption{On-Device Hardware Profiling Across Edge Engine Catalog}
    \label{tab:mobile_telemetry}
    \renewcommand{\arraystretch}{1.2}
    \resizebox{\textwidth}{!}{
    \begin{tabular}{lcccc}
    \toprule
    \textbf{Engine Spec} & \textbf{File Size (GGUF)} & \textbf{RAM Footprint} & \textbf{Inference Speed} & \textbf{Cold-Start Latency} \\
    \midrule
    Qwen2.5-0.5B-Instruct (Q4\_K\_M) & 392 MB & 580 MB & 41.2 tok/s & 180 ms \\
    Llama-3.2-1B-Instruct (Q4\_K\_M) & 804 MB & 1,040 MB & 28.6 tok/s & 320 ms \\
    Gemma-2-2B-Instruct (Q4\_K\_M) & 1,629 MB & 2,200 MB & 22.4 tok/s & 640 ms \\
    Llama-3.2-3B-Instruct (Q4\_K\_M) & 1,925 MB & 2,600 MB & 18.2 tok/s & 850 ms \\
    \bottomrule
    \end{tabular}
    }
    \end{table}

    \section{Deterministic Tool Execution and Intent Routing}
    The agent writes four primary autonomous operational pillars:
    \begin{enumerate}
        \item \textbf{Local Reminders (\texttt{create\_reminder}):} Schedules exact Android alarm notifications stored within the encrypted Room SQLite tables.
        \item \textbf{Calendar Search (\texttt{search\_calendar}):} Queries local calendar events by date and keyword.
        \item \textbf{Private Notes Management (\texttt{manage\_notes}):} Creates, searches and reads encrypted personal notes.
        \item \textbf{Application Dispatcher (\texttt{open\_app}):} Solves and launches permitted installed packages (WhatsApp, Spotify, Google Maps, Camera, Calculator, Settings), in real-world scenarios.
    \end{enumerate}

    To make sure complete privacy and user safety, plain two-phase human authorization is generally needed before intent dispatch by actions with external side effects (e.g., \texttt{draft\_message}), on the whole.

    % ==============================================================================
    % CHAPTER 7: BIT-BUDGET GUIDELINES
    % ==============================================================================
    \chapter{Practical Bit-Budget Guidelines for Edge Engineers}

    \section{The 3-Percentage-Point ($\le 3\,\text{pp}$) Acceptability Threshold}
    In edge systems engineering, an tuning can't be judged solely by parameter compression ratio; it must satisfy minimum task fidelity. We define a rollout arrangement as \textbf{Acceptable} if and only if:
    \begin{equation}
        \Delta\text{Acc}_{\text{task}} = \text{Acc}_{\text{FP16}} - \text{Acc}_{\text{quant}} \le 3.0\,\text{percentage points}
    \end{equation}

    \section{Task-Conditioned Bit-Budget Decision Matrix}
    Based on our 30-run empirical evaluation grid, doable engineering guidelines for choosing quantization bit-widths across the separate application domains are being offered by Table~\ref{tab:bit_guidelines}, on the whole.

    \begin{table}[h!]
    \centering
    \caption{Task-Conditioned Bit-Budget Decision Matrix for Edge Computing}
    \label{tab:bit_guidelines}
    \renewcommand{\arraystretch}{1.2}
    \resizebox{\textwidth}{!}{
    \begin{tabular}{p{3.8cm}cp{4.2cm}p{4cm}}
    \toprule
    \textbf{Workload Domain} & \textbf{Min Bit-Budget} & \textbf{Recommended Quant Format} & \textbf{Engineering Risk Factors} \\
    \midrule
    Factual QA \& Static Retrieval & 4-bit & GPTQ / AWQ / Q4\_K\_M & Minimal degradation ($\le 1.5\,\text{pp}$). Safe for edge deployment. \\
    Single-Step Classification & 4-bit & Standard RTN / GPTQ-4bit & Highly resilient. Even 3-bit remains viable ($\le 3.5\,\text{pp}$). \\
    Structured Grammar Tool Calling & 4-bit & GGUF Q4\_K\_M + GBNF & High syntax sensitivity. Requires GBNF logit masking. \\
    Multi-Step Symbolic Math & \textbf{8-bit} & GPTQ-Int8 / Q8\_0 & \textbf{High Risk at 4-bit} ($>4\,\text{pp}$ drop). 2-bit collapses completely. \\
    Code Generation & 8-bit & AWQ-Int8 / GGUF Q8\_0 & Indentation and syntax collapse under aggressive quantization. \\
    \bottomrule
    \end{tabular}
    }
    \end{table}

    % ==============================================================================
    % CHAPTER 8: PROJECT TIMELINE & GANTT CHART
    % ==============================================================================
    \chapter{Work Plan, Milestones, and Gantt Schedule}

    \section{Semester Project Schedule and Deliverables}
    During the 2026--2027 academic session, the project was ran across a 10-week timeline:
    \begin{itemize}[leftmargin=*]
        \item \textbf{Weeks 1--2 (Environment Setup):} Pipeline configuration, pinned \texttt{lm-evaluation-harness} validation, and Kaggle T4 environment setup.
        \item \textbf{Weeks 3--4 (Baseline Runs):} Execution of unquantized FP16 baselines across Qwen2.5, Llama-3.2, and Pythia on MMLU, ARC and HellaSwag, in real-world scenarios.
        \item \textbf{Weeks 5--6 (Full Quantization Grid):} Quantization of models into GPTQ, AWQ and NF4 formats; execution of GSM8K and IFEval benchmarks, on the whole.
        \item \textbf{Weeks 7--8 (Statistical Analysis):} Computation of Spearman rank correlations, bootstrap confidence intervals, and perplexity reliability analysis.
        \item \textbf{Weeks 9--10 (Edge Agent Deployment & Thesis):} Android native bridge development, ARM NEON profiling, and final dissertation preparation.
    \end{itemize}

    \begin{table}[h!]
    \centering
    \caption{10-Week Minor Project Execution Milestones}
    \label{tab:gantt_table}
    \renewcommand{\arraystretch}{1.2}
    \resizebox{\textwidth}{!}{
    \begin{tabular}{clll}
    \toprule
    \textbf{Week} & \textbf{Milestone Description} & \textbf{Primary Deliverable} & \textbf{Status} \\
    \midrule
    W1--W2 & Environment Setup \& Dry Runs & Pinned harness repo, Kaggle script pipelines & Completed \\
    W3--W4 & FP16 Baseline Acquisition & Verified FP16 accuracy matching model cards & Completed \\
    W5--W6 & Full Quantization Matrix Runs & Master table of 27 quantized configurations & Completed \\
    W7--W8 & Statistical Analysis \& Spearman Correlation & $\rho$ matrix, perplexity decoupling proof & Completed \\
    W9--W10 & Mobile Agent Deployment \& Final Viva & Production APK v1.0.10, Final LaTeX Report & Completed \\
    \bottomrule
    \end{tabular}
    }
    \end{table}

    % ==============================================================================
    % CHAPTER 9: CONCLUSION & FUTURE WORK
    % ==============================================================================
    \chapter{Conclusion, Limitations, and Future Work}

    \section{Summary of Research Contributions}
    This B.Tech Minor Project delivers three primary advancements to the discipline of machine learning systems:
    \begin{enumerate}
        \item \textbf{Rigorous Empirical Benchmarking:} A controlled check across the compact LLM architectures was did, creating an complete 30-run comparison table across six downstream tasks, in many respects.
        \item \textbf{Deconstruction of the Perplexity Fallacy:} We proved that language-modelling perplexity decouples from multi-step reasoning accuracy ($\rho = 0.41$ on GSM8K), warning researchers against counting on perplexity for edge deployment decisions, on the whole.
        \item \textbf{Production On-Device Agent Deployment:} Empirical ideas were translated into a functional Android mobile agent running Google Gemma-2-2B and Llama-3.2 locally employing Q4\_K\_M GGUF quantization, ARM NEON SIMD optimization, and GBNF grammar constraints, reaching true zero-cloud-egress autonomous device assistance, broadly speaking.
    \end{enumerate}

    \section{Limitations}
    Our check is often subject to two primary shortcomings:
    \begin{itemize}
        \item \textbf{Hardware Scope:} GPU benchmark evaluations were did mostly on NVIDIA T4 architectures, really, even as mobile profiling was inspected on Snapdragon 8 Gen 2, broader cross-mobile SoC testing, across MediaTek Dimensity and Google Tensor is often remaining ongoing, broadly speaking.
        \item \textbf{Activation Quantization:} What's more our core research stressed on weight-only post-training quantization ($W_q A_{16}$); thus, joint weight-activation quantization ($W_8 A_8$ or $W_4 A_4$) introduces activation outlier clipping that warrants further exploration, broadly speaking.
    \end{itemize}

    \section{Future Directions}
    Future extensions cover:
    \begin{itemize}
        \item Combining KV cache quantization (K-quants 4-bit KV cache) to support 32k context windows on mobile devices, to a significant degree.
        \item Studying speculative decoding on mobile CPUs employing ultra-compact 0.5B draft models paired with 3B verification engines, on the whole.
    \end{itemize}

    % ==============================================================================
    % REFERENCES
    % ==============================================================================
    \begin{thebibliography}{99}
    \addcontentsline{toc}{chapter}{Bibliography}

    \bibitem{frantar2023gptq}
    E.~Frantar, S.~Ashkboos, T.~Hoefler, and D.~Alistarh,
    ``GPTQ: Accurate post-training quantization for generative pre-trained transformers,''
    in \textit{Proc. International Conference on Learning Representations (ICLR)}, 2023.

    \bibitem{lin2024awq}
    J.~Lin, J.~Tang, H.~Tang, et~al.,
    ``AWQ: Activation-aware weight quantization for on-device LLM compression and acceleration,''
    in \textit{Proc. Machine Learning and Systems (MLSys)}, 2024.

    \bibitem{dettmers2023qlora}
    T.~Dettmers, A.~Pagnoni, A.~Holtzman, and L.~Zettlemoyer,
    ``QLoRA: Efficient finetuning of quantized LLMs,''
    in \textit{Advances in Neural Information Processing Systems (NeurIPS)}, vol.~36, 2023.

    \bibitem{dettmers2022llmint8}
    T.~Dettmers, M.~Lewis, Y.~Belkada, and L.~Zettlemoyer,
    ``LLM.int8(): 8-bit matrix multiplication for transformers at scale,''
    in \textit{Advances in Neural Information Processing Systems (NeurIPS)}, vol.~35, 2022.

    \bibitem{jin2024comprehensive}
    R.~Jin, J.~Du, W.~Huang, et~al.,
    ``A comprehensive evaluation of quantization strategies for large language models,''
    in \textit{Findings of the Association for Computational Linguistics: ACL 2024}, pp.~3210--3229, 2024.

    \bibitem{li2024evaluating}
    S.~Li, X.~Ning, L.~Wang, et~al.,
    ``Evaluating quantized large language models,''
    in \textit{Proc. International Conference on Machine Learning (ICML)}, 2024.

    \bibitem{chee2023quip}
    J.~Chee, N.~Cai, V.~Kuleshov, and C.~De~Sa,
    ``QuIP: 2-bit quantization of large language models with incoherence processing,''
    in \textit{Advances in Neural Information Processing Systems (NeurIPS)}, vol.~36, 2023.

    \bibitem{tseng2024quip}
    A.~Tseng, J.~Chee, Q.~Sun, V.~Kuleshov, and C.~De~Sa,
    ``QuIP\#: Even better LLM quantization with Hadamard transform and vector codebooks,''
    in \textit{Proc. International Conference on Machine Learning (ICML)}, 2024.

    \bibitem{dong2020hawq}
    Z.~Dong, Z.~Yao, D.~Arfeen, A.~Gholami, and M.~W. Mahoney,
    ``HAWQ-V2: Hessian aware trace-weighted quantization of neural networks,''
    in \textit{Advances in Neural Information Processing Systems (NeurIPS)}, 2020.

    \bibitem{yao2021hawq}
    Z.~Yao, Z.~Dong, Z.~Zheng, A.~Gholami, et~al.,
    ``HAWQ-V3: Dyadic neural network quantization,''
    in \textit{Proc. International Conference on Machine Learning (ICML)}, 2021.

    \bibitem{kim2023squeezellm}
    S.~Kim, C.~Hoover, E.~Gholami, et~al.,
    ``SqueezeLLM: Dense-and-sparse quantization,''
    in \textit{Proc. International Conference on Machine Learning (ICML)}, 2023.

    \bibitem{egorov2024extreme}
    Y.~Egorov, A.~Tseng, et~al.,
    ``AQLM: Extreme compression of large language models with additive quantization,''
    in \textit{Proc. International Conference on Machine Learning (ICML)}, 2024.

    \bibitem{guo2024vptq}
    Y.~Guo, Z.~Zhao, et~al.,
    ``VPTQ: Extreme low-bit vector post-training quantization for large language models,''
    \textit{arXiv preprint arXiv:2409.17066}, 2024.

    \bibitem{qwen2024}
    Qwen Team,
    ``Qwen2.5: A party of foundation models,''
    \textit{arXiv preprint arXiv:2412.15115}, 2024.
    
    \bibitem{dubey2024llama}
    A.~Dubey, A.~Jauhri, A.~Pandey, et~al.,
    ``The Llama 3 herd of models,''
    \textit{arXiv preprint arXiv:2407.21783}, 2024.

    \bibitem{biderman2023pythia}
    S.~Biderman, H.~Schoelkopf, Q.~Anthony, et~al.,
    ``Pythia: A suite for analyzing large language models across training and scaling,''
    in \textit{Proc. International Conference on Machine Learning (ICML)}, pp.~2397--2430, 2023.

    \bibitem{team2024gemma}
    Gemma Team,
    ``Gemma 2: Improving open language models at a practical size,''
    \textit{arXiv preprint arXiv:2408.00118}, 2024.

    \end{thebibliography}

    \end{document}